\chapter{Βάση δεδομένων}

\section{Σχεσιακό σχήμα σε λογικό επίπεδο}
\label{sec:schema-logical}

Η σχεδίαση ακολουθεί κλασικό \emph{Kimball star schema} με \emph{wide dimensions}.
Οι λόγοι της επιλογής αυτής:

\begin{itemize}
    \item \textbf{Single-version-of-the-truth lookup tables} (\texttt{lookup\_authors},
          \texttt{lookup\_journals}, \texttt{lookup\_conferences},
          \texttt{lookup\_field\_of\_research\_categories}) με ακέραια \texttt{INT
          UNSIGNED} primary keys --- οι συμβολοσειρές δεν αποτελούν αποδεκτά PKs
          (όπως ρητά αναφέρει η εκφώνηση).
    \item \textbf{Fact tables} (\texttt{fact\_journal\_articles},
          \texttt{fact\_conference\_articles}) με τα ίδια PKs (article\_id) και
          FKs στα lookups· διαφορετική δομή ανάλογα με τον τύπο venue, όπως ζητείται.
    \item \textbf{Bridge tables} (\texttt{bridge\_journal\_article\_authors},
          \texttt{bridge\_conference\_article\_authors}) για την N:M σχέση
          άρθρων--συγγραφέων με σύνθετο PK.
    \item \textbf{Quarantine table} (\texttt{rejection\_logs}) που καταγράφει
          αναλυτικά οποιαδήποτε γραμμή απορρίφθηκε από τον ETL transformer ---
          δεν χάνεται καμία πληροφορία, ούτε σιωπηλά απορρίπτεται κάτι.
    \item \textbf{Wide journal dimension}: όλα τα Kaggle Scimago πεδία
          (\texttt{sjr\_index}, \texttt{cite\_score}, \texttt{h\_index},
          \texttt{total\_docs}, \texttt{total\_cites\_3y} κ.λπ.) ζουν μέσα στο
          \texttt{lookup\_journals} αντί για ξεχωριστό
          \texttt{journal\_ranking} table. Trade-off: ευκολότερα reads, ελάχιστα
          extra bytes ανά γραμμή (περίπου 120 B), χωρίς denormalisation overhead από
          άποψη ενημερώσεων (τα ranking δεδομένα είναι read-only μετά το ETL).
\end{itemize}

\paragraph{Διάγραμμα.} Το αναλυτικό ERD των πινάκων δίνεται στο
Σχήμα~\ref{fig:schema-tables}, και η αντίστοιχη όψη των aggregate views
που τους τροφοδοτούν στο Σχήμα~\ref{fig:schema-views}.

\begin{figure}[H]
    \centering
    \includegraphics[width=\linewidth,height=0.78\textheight,keepaspectratio]{erd.png}
    \caption{Σχεσιακό σχήμα: lookup, fact, bridge και materialized tables
    (reverse-engineered στο DataGrip). Τα βέλη δείχνουν τα foreign keys, οι
    materialized πίνακες είναι read-only snapshots που γεμίζουν στο τέλος
    του \texttt{03\_views.sql}.}
    \label{fig:schema-tables}
\end{figure}

\begin{figure}[H]
    \centering
    \includegraphics[width=\linewidth,height=0.78\textheight,keepaspectratio]{views.png}
    \caption{SQL views που τρέχουν πάνω από το σχεσιακό σχήμα. Τα
    \texttt{view\_*\_profile} εξυπηρετούν τις σελίδες προφίλ, τα
    \texttt{view\_*\_yearly\_statistics} τα line charts, και το
    \texttt{view\_corpus\_totals} τα KPI tiles του dashboard.}
    \label{fig:schema-views}
\end{figure}

\section{Σχεσιακό σχήμα σε φυσικό επίπεδο}

\subsection{Ρύθμιση παραμέτρων DBMS}

\begin{itemize}
    \item \textbf{DBMS}: MySQL 8.0 σε Docker container (image
          \texttt{mysql:8.0}). Λόγος: ελαφρύ, νέα έκδοση με υποστήριξη
          \texttt{WITH} statements (CTEs), εξαιρετική υποστήριξη
          \texttt{LOAD DATA LOCAL INFILE}.
    \item \textbf{Storage engine}: \texttt{InnoDB} σε όλους τους πίνακες.
    \item \textbf{Charset / collation}: \texttt{utf8mb4} /
          \texttt{utf8mb4\_unicode\_ci} --- απαραίτητο για διεθνή
          ονόματα συγγραφέων (Asian + Cyrillic + Greek).
    \item \textbf{Container flags}: \texttt{-{}-local-infile=1} ώστε να μπορεί
          να εκτελεστεί \texttt{LOAD DATA LOCAL INFILE} από τον client.
    \item \textbf{Connection pool} (asyncmy): min 1, max 8 connections,
          recycle 1800s.
\end{itemize}

\subsection{Ρύθμιση φυσικού σχήματος}
\label{sec:physical}

\textbf{Indexes.} Πέρα από τα PKs/FKs, ορίσαμε σύνθετα indexes σε:

\begin{itemize}
    \item \texttt{fact\_journal\_articles (journal\_id, year)} και αντίστοιχα
          για conference articles, ώστε τα filtered year-range queries να
          τρέχουν με index-range scan αντί full table scan.
    \item \texttt{lookup\_journals (best\_subject\_area)} και
          \texttt{lookup\_journals (publisher)} για τα filter dropdowns.
    \item \texttt{materialized\_author\_profile (total\_articles)} και
          \texttt{(total\_articles, author\_name)} για τον \texttt{ORDER BY}
          του authors browse table.
\end{itemize}

\textbf{Views} (11 συνολικά) ορίζονται στο \texttt{sql\_scripts/03\_views.sql}:
\texttt{view\_journal\_profile}, \texttt{view\_journal\_yearly\_statistics},
\texttt{view\_conference\_profile}, \texttt{view\_conference\_yearly\_statistics},
\texttt{view\_year\_summary}, \texttt{view\_corpus\_totals},
\texttt{view\_publisher\_quartile\_distribution},
\texttt{view\_subject\_area\_yearly\_summary},
\texttt{view\_field\_of\_research\_yearly\_summary},
\texttt{view\_author\_profile}, \texttt{view\_author\_yearly\_statistics}.

\textbf{Materialised tables} (post-ETL warehouse pattern --- οι αναφορές διαρκούν
σταθερά λιγότερο από 50 ms αντί δευτερολέπτων):

\begin{itemize}
    \item \texttt{materialized\_year\_summary}
    \item \texttt{materialized\_subject\_area\_yearly\_summary}
    \item \texttt{materialized\_field\_of\_research\_yearly\_summary}
    \item \texttt{materialized\_authors\_vs\_articles\_scatter\_journals}
    \item \texttt{materialized\_authors\_vs\_articles\_scatter\_conferences}
    \item \texttt{materialized\_author\_profile}
\end{itemize}

Οι materialised πίνακες παράγονται από τα αντίστοιχα views μέσω
\texttt{CREATE TABLE \dots AS SELECT} στο τέλος του \texttt{03\_views.sql},
μετά τη φόρτωση των fact tables.

\subsection{Ρύθμιση ασφάλειας}

Στη φάση παράδοσης, το σύστημα τρέχει με ένα μόνο role (root) μέσα σε
container. Οι credentials δεν είναι hardcoded πουθενά --- φορτώνονται από
\texttt{.env} (\texttt{MYSQL\_ROOT\_PASSWORD}, \texttt{MYSQL\_DATABASE},
\texttt{MYSQL\_USER}, \texttt{MYSQL\_PASSWORD}) μέσω
\texttt{pydantic-settings}.

\paragraph{Παραμετροποιημένα queries.} Ολα τα queries του back-end χρησιμοποιούν
\texttt{cur.execute(\dots WHERE id = \%s, (id,))} --- δεν υπάρχει f-string
ή \texttt{format()} interpolation σε SQL string πουθενά στον κώδικα.
Αυτό αποκλείει SQL injection ακόμα και αν ο front-end υποστεί compromise.
